% 第六部：結果面的理由（第肆部分，1376–1386）。出處：細目次第三節肆；W06_1019_導讀問答.md 案件表與理論表（已查核）；英文底稿。
\chapter{結果面的理由：法院是否較能得出正確的權利決定}

\qaentry{第六部之一}{第肆部分的出發點}

\begin{qaquestion}
第肆部分從誰的主張出發？作者為何說結果面的理由對兩個機關都有？
\end{qaquestion}

第肆部分從 Raz 的出發點與法律程序學派關於機構能力的論辯起論\pg{1376}。Raz 指控立法機關受宗派利益驅使；作者答以兩點。其一，同樣的宗派壓力也說明法院對權利的忽視，正文舉 Korematsu、Schenck、Dred Scott、Prigg 等案（註 82），並引通常堅定支持違憲審查的 Tribe 在九一一之後的評論：危機時期不宜對司法寄予過高的期望\pg{1377}。其二，立法機關若普遍如此，即與第三項假設不相容，屬非核心情形\pg{1377}。Raz 自己也承認結果面的理由有相反的方向：受決定影響的人較能體察處境，立法機關的代表結構正為培養這種體察\pg{1377--1378}。

民主程序本身已含有為結果而作的調整，例如選舉權的限制與立法機關的組成，屬 Rawls 所稱的不完全的程序正義\pg{1378}。作者接著指出，違憲審查的支持者對司法決定的描述不切實際\pg{1379}，以下逐一檢驗被歸給法院的三項優點。

\qaentry{第六部之二}{被歸給法院的三項優點}

\begin{qaquestion}
支持者認為法院在哪三方面較能得出正確的權利決定？\\
作者如何逐一回應？
\end{qaquestion}

\begin{longtable}{@{}L{0.16\linewidth}L{0.30\linewidth}L{0.46\linewidth}@{}}
\toprule
優點 & 支持者的主張 & 作者的回應 \\
\midrule
\endhead
A. 以個案為取向（1379--1380） & Moore：法官每天面對具體個人的處境，較能得到道德洞見 & 「But this is mostly a myth.」案件上訴到最高審級時，原權利人的痕跡幾已消失，論辯圍繞抽象的權利問題；原告常由倡議團體挑選\pg{1379--1380}。立法程序同樣考量個案，Megan's Law 即因知名個案而立法；「困難案件造就惡法」若屬實，立法機關較能評估個案與影響數百萬人的一般權利問題之間的關係\pg{1380} \\
B. 以權利法案的文本為取向（1380--1382） & 法院依附權利法案的條文推理 & 條文的措辭往往為迴避制定時的分歧而選；解釋理論的論辯排擠對道德問題的直接論辯；法院因自身正當性的疑慮而依附文本；文本所遺漏的權利（如社會經濟權利）同樣扭曲推理 \\
C. 說明理由（1382--1386） & 法院必須公開說明理由 & 議員同樣說明理由，載於 Hansard 與 Congressional Record；法院的理由以條文、先例為中心，成因是法院須證明自己有權作成該決定 \\
\bottomrule
\end{longtable}

\qaentry{第六部之三}{1966 年英國下議院墮胎法案辯論與 Roe 案}

\begin{qaquestion}
作者為何拿英國國會的一場辯論與 Roe v.\ Wade 判決對比？對比的內容為何？
\end{qaquestion}

作者寫到，在權利的善意分歧中真正的爭點常被推到邊緣：在 Roe v.\ Wade 五十頁的意見書中，處理生育權與隱私之關係的道德重要性只有幾段，處理胎兒地位的段落多半在呈現各種意見的分歧\pg{1383}；註中說該意見書含大量的法律史與社會史。

對照的是 1966 年英國下議院 Medical Termination of Pregnancy Bill 的二讀辯論。作者讀過議事錄，約一百頁的辯論直接處理墮胎的道德問題，連反對陣營的議員也稱許其立場受到尊重的傾聽\pg{1384--1385}；註 106 舉投票反對該法案的天主教議員 Norman St.\ John-Stevas，稱辯論的水準無愧於下議院最高的傳統。作者接著問，反對墮胎的一方何曾稱許最高法院在 Roe 案中對其立場的關注與尊重\pg{1385}。

\qaentry{第六部之四}{原則的論壇與常態、病態之分}

\begin{qaquestion}
作者如何回應 Dworkin「法院是原則的論壇」之說？\\
立法機關的缺陷為何不計入比較？
\end{qaquestion}

對 Dworkin 的回應，作者指出法院的推理以解釋與教義為主，英國的立法辯論能直接處理道德問題\pg{1385--1386}。他承認立法的推理也可能是 a disgrace：

\begin{quote}
\ldots so too can legislative reasoning be a disgrace, as legislative majorities act out of panic, recklessly, or simply parrot popular or sectarian slogans in their pseudo-debates.\pg{1386}
\end{quote}

他接著問兩個機關的缺陷何者屬常態、何者屬偏差。法院依先例、文本、教義推理正是被期待的行為方式，第二項假設所描述的就是依這些標準運作良好的法院；倉促或宗派性的立法則不屬立法機關被設立來做的事，不應假定於核心情形\pg{1386}。

\begin{readernote}
法院的缺陷計入常態，在 B、C 兩節拿來與立法機關比較；立法機關的缺陷劃為病態，移到第柒部分。這項劃分的依據是兩個機關設立的目的，作者未就兩者的實際運作作經驗上的比較。若以設立目的為準，法院的設立目的同樣可以說成就權利問題說理，法條主義也就可以歸為偏差。討論時可從這一段提問。
\end{readernote}
